\begingroup
\small
\setlength{\tabcolsep}{4pt}
\renewcommand{\arraystretch}{1.12}
\begin{longtable}{|L{0.25\textwidth}|L{0.63\textwidth}|}
\caption{Нефункциональные требования}\label{tab:nonfunctional-requirements}\\
\hline
\textbf{Параметр} & \textbf{Принятое условие} \\
\hline
\endfirsthead
\multicolumn{2}{c}{Продолжение таблицы~\thetable}\\
\hline
\textbf{Параметр} & \textbf{Принятое условие} \\
\hline
\endhead
Размер страницы выдачи & 20 аудиозаписей на странице с пагинацией. \\
\hline
Загрузка аудиофайлов & Загрузка аудиофайлов происходит на стороне сервиса каталога, подсистема анализа и поиска лишь получает доступ к уже сохранённым записям, обрабатывает их и хранит результаты анализа. \\
\hline
Хранение аудиофайлов & Аудиофайлы хранятся в S3-хранилище. \\
\hline
Внешние зависимости & Проверка прав через сервис аутентификации занимает не более 200 мс; скорость чтения аудиозаписей из S3-хранилища - не менее 10 Мб/с. \\
\hline
Поддерживаемые аудиозаписи & Аудиоайлы с расширениями MP3, WAV и FLAC, размером до 100 Мб и длительностью до 15 минут. \\
\hline
Дизайн интерфейса & Интерфейс должен быть удобен, интуитивно понятен и выполнен в единой дизайн-системе. \\
\hline
Качество распознавания эмоций & Значение macro-F1 - не ниже \textbf{0,75} на независимой размеченной выборке из не менее чем 500 аудиозаписей. Для каждой метки нужны минимум 50 положительных и 50 отрицательных примеров. \\
\hline
Уверенность модели & Эмоциональная метка принимается, если оценка уверенности модели не ниже \textbf{0,7}; меньшая оценка считается недостаточной. \\
\hline
Длительность анализа & Для большинства файлов анализ должен завершаться за $2 \times t$ секунд, где $t$ - длительность аудиозаписи. Время ожидания в очереди измеряется отдельно. \\
\hline
Нагрузка & Сервис должен выдерживать 10 одновременных клиентов и 10 запросов в секунду. \\
\hline
Актуальность поиска & Результат, пригодный для поиска, появляется в индексе не позднее чем через \textbf{30 секунд} после сохранения. При массовом переанализе поиск продолжает работать по текущей версии, пока новая версия индекса не будет готова и активирована. \\
\hline
Количество сбоев & Допускается не более одного сбоя за календарный квартал. \\
\hline
Восстановление после сбоя & Среднее время восстановления сервиса после сбоя должно составлять не более чем 1,5 часа. \\
\hline
Сохранность заданий анализа & При перезапуске обработчика принятые задания и сохранённые результаты не должны теряться. \\
\hline
Резервное восстановление & После повреждения хранилища допускается потерять не более 30 минут изменений. Восстановление должно занимать не более 3 часов; его проверяют минимум раз в квартал. \\
\hline
Согласованность данных & Результат сохраняется вместе с версиями аудиозаписи, модели и настроек. Незавершённый результат не показывается пользователю и не попадает в поиск. \\
\hline
Замена модели & Новую совместимую модель можно подключить через отдельный адаптер, не меняя внешний API и основные сценарии сервиса. \\
\hline
Масштабируемость & Дополнительные обработчики аудио можно запускать без изменения внешнего API. Совместная работа обработчиков не должна приводить к потере заданий или результатов. \\
\hline
Сохранность оригиналов аудиозаписей & Подсистема имеет доступ только на чтение аудиофайлов. \\
\hline
\end{longtable}
\endgroup
